\section{Introduction}

Consumer-finance workflows increasingly route individual decisions through
large language models (LLMs): a support message is routed to hardship review,
a flagged transaction is cleared or escalated. The attributes consulted in
these decisions are purpose-bound by law and by contract. Under the principle
of purpose limitation and data minimization~\cite{gdpr2016,
purposeminimization2021}, a lender may hold a hardship flag to evaluate
repayment-assistance eligibility (authorized purpose) while being prohibited
from using the same flag to decide service priority (prohibited purpose). The
European AI Act similarly treats credit and insurance scoring as high-risk
applications with strict purpose boundaries~\cite{euaiact2024}.

Modern access-control and privacy machinery addresses where data may \emph{go}
(transmission), not what it may \emph{do} at the decision. Confidential
computing~\cite{opaque2017,chiron2018,slalom2019}, usage
control~\cite{usagecontrol2018}, and purpose-based access control
(PBAC)~\cite{byun2008} govern release. Prompt-injection defenses~\cite{agentdojo2024,camel2025}
protect the prompt. But a model that is allowed to \emph{receive} a field for an
authorized purpose can still \emph{act on} it under a prohibited purpose: the
field is visible in the prompt, the model weighs it, and the decision drifts.
Transmission-level controls neither measure nor prevent this drift.

We make three contributions.

\begin{enumerate}
  \item \textbf{Definition.} We formalize \emph{Purpose-Selective Bounded
  Execution (PSBE)} as two measurable decision-level properties over a
  confidential attribute: the \emph{authorized utility} it preserves for the
  permitted purpose, and the \emph{prohibited influence} it exerts under other
  purposes, measured relative to the system's own natural-decision floor.
  \item \textbf{Benchmark.} We introduce \emph{FinBoundBench}, a paired
  benchmark of synthetic-but-realistic consumer-finance cases in which the same
  confidential field (e.g., a hardship flag, a fraud flag) is embedded in
  purpose-labeled decision tasks. Each case is rendered into counterfactual
  variants (authorized, prohibited, masked), so a single model serves as its
  own control.
  \item \textbf{Evidence.} We report a preregistered, frozen, and independently
  re-computed confirmatory study across two model families (DeepSeek V4 Pro;
  Kimi K3) and two financial tasks (hardship-support routing; fraud review).
  Prohibited-visible data changed decisions on 96--100\% of cases (floor:
  0--4\%); masking the field suppressed influence to at or below the floor in
  every condition; and governed execution with a declared purpose retained the
  full authorized utility gain (utility-retention ratio of exactly 1.0 in both
  studies).
\end{enumerate}

Our framing is deliberately modest. The study covers two model-task pairs over
semi-synthetic signals; it does not claim formal guarantees, trusted-hardware
isolation, or adversarial robustness (Section~\ref{sec:limitations}). Its
contribution is a measurement methodology that isolates the quantity that
matters---\emph{decision influence relative to a floor}---together with
confirmatory evidence that a governed execution adapter can be purpose-selective
without an authorized-utility penalty.
